\section{Realizing $\F_\oslone$ with OT with Less Precision}
\label{apx:otl1}
We implement oblivious sampling using a 1-out-of-$m$ OT scheme. In
particular, the receiver, as an OT receiver, chooses a random index from
$[m]$, and the sender, as an OT sender, prepares an $m$-dimensional
input vector that encodes the $L_1$ distribution of $\vec{w}$ in a way
that we will describe soon.

\paragraph{Assumption about the level of precision of the input.} 
With this approach, each element from the prepared OT input vector will
be chosen uniformly with probability $\frac 1 m$. Therefore, the size
$m$ affects the level of precision of the sampling.
%
In particular, we set $\mu := 1/m$ as a precision unit, and we assume
the following:

\BI
\item[]
  \it For each $i \in [n]$, it holds that $\frac{w_{i}}{\|\vec{w}\|_1}$ is a multiple of $\mu$.
\EI
%
If the input vector $\vec w$ is not consistent with the above
requirement, one can round it by using the following function
$\rounding_{\mu}(\vec{w})$: 

\BEN
\item[] \hspace{-2em} $\rounding_\mu(\vec{w})$
\item Let $\vec{w} = (w_1, \ldots, w_n)$.  For $i = 1, \ldots, n$,
compute $w_i' = \trunc_{\mu}(w_i)$. Here, for any real number $x$ with $x \in [0,1]$, denote $\trunc_{\mu}(x) =
\tilde x \cdot \mu$ where $\tilde x$ is an integer that minimizes $\Delta
:= x - \tilde x \cdot \u$ subject to $\Delta \ge 0$. 
Typically, we have $\mu = 2^{-q}$
for a certain positive integer $q$, and $\trunc_{\mu}(x)$ is simply
truncating the lower order bits in the binary representation of $x$. 



Let $\vec{w}' = (w_1', \ldots,
w_n')$. 

\item Repeat the following until $L_1$ norm of $\vec{w}'$ becomes 1. 
    \BI
    \item[] Find $j = \arg\max_{i \in [n]} (w_i - w_i')$, and increase
    $w'_j$ by $\mu$.
    \EI
\item Output $\vec{w}'$. 
\EEN

The above algorithm makes sure that for all $i$, it holds $|w'_i - w_i|
< \mu$, which means $\vec{w}'$ is a good approximation of $\vec{w}$, with
each element having an additive error of at most $\mu$. To see why, note
that in step 1, some $w'_i$s will get truncated leading to small
difference, i.e., $w_i - w_i' < \mu$. In step 2, since the truncated
weights are added back to the elements in decreasing order of difference
$(w_i - w_i')$, only some of the truncated $w'_i$s will be updated to
$w'_i + \mu$ (which will still be close to $w_i$) until the $L_1$ norm
of $\vec{w}'$ becomes 1. 

In a situation where the low precision is acceptable, this OT-based
solutions could be more efficient. However, if one needs a higher level
of precision, we recommend using the FHE-based solution described in the
next subsection.  We also observe that by using an OT protocol with
$O(\log m)$ communication~\cite[Theorem 2.2]{IMS+09}, we expect we could
support fairly large values of $m$.  The bottleneck on larger $m$ is
likely to be storage, and the computation time needed for the OT.  

\floatname{algorithm}{Protocol}
\begin{algorithm}[htbp]
\textbf{Inputs:} The sender has input $\vec{w}$. We require every
$w_{i}/\|\vec{w}\|_1$ is a multiple of $\mu := \frac 1 m$. 
\BEN
\item The sender computes the following: 

  \BEN
  \item 
    Given $\vec{w}$, the sender prepares an $m$-dimensional input vector as follows: 

  \item 
   For $i = 1, \ldots, n$, do: 
      \BEN
      \item[] Let $k_i = \frac{w_i}{\|\vec{w}\|_1} \cdot m$.  Insert $k_i$ copies of the index
      $i$ into the $m$-dimensional  vector $\vec{v}$; that
      is, there should be $k_i$ slots (out of $m$)  whose value is $i$ in the
      $m$-dimensional vector $\vec{v}$. 
      \EEN
    Note that for each $i$, the fraction of the slots containing the index
    $i$ in $\vec{v}$ is $\frac{k_i}{m} = \frac{w_i}{\|\vec{w}\|_1}$.  

  \item The sender chooses a random pad $\pi$ uniformly at random.

  \item Let $\vec{v} = (v_{1}, \ldots, v_{m})$. The sender shuffles
    $\vec{v}$ and blinds it by updating $v_{i} := v_{i} \xor \pi$ with a
    randomly chosen $\pi$.
  \EEN

\item Execute an OT protocol where the sender is the OT sender with input
$\vec{v}$ and the receiver is the OT receiver with a randomly chosen number from
$[m]$. Let $u$ be the output to the OT receiver.

\item Output $\pi$ to the sender and output $u$ to the receiver. 
\EEN

\caption{Oblivious sampling protocol realizing $\F_{\oslone}$ based on
$1$-out-of-$m$ OT}
\label{fig:osample-ot}
\end{algorithm}


\paragraph{Oblivious sampling protocol.} 
With the assumption about the level of precision of the input vector
$\vec{w}$, we can implement oblivious sampling. See
Protocol~\ref{fig:osample-ot}. 

Due to security of the OT protocol, the OT sender won't know the OT
receiver's choice. However, since the OT receiver does know the sampled
index, which leaks information about the data array, we hide this
information from the OT receiver by having the OT sender {\em shuffle}
the input vector.     

At the end of the OT protocol, the sender and receiver will hold the
sampled index $i$ in a secret shared form; that is, the sender will hold
$\pi$ and the receiver $\pi \xor i$.  Note that the sender re-uses the
same $\pi$ across all inputs, in order to fix its share, independently
of the receiver index.  Security holds even with the re-use of this
value because the receiver learns only a single element.


\paragraph{Security.}
We will prove the following theorem. 
\begin{theorem} \label{thm:osample-ot}
Protocol~\ref{fig:osample-ot} securely realizes $\F_{\oslone}$ in the
semi-honest security model. 
\end{theorem} 


\begin{proof}
We first give the simulation of the sender. This is trivial in the OT-hybrid model, as the sender receives no messages in this protocol.  The simulator submits the sender's input to $\F_{\oslone}$, and recieves $\pi$ as output.  It places $\pi$ on the sender's random tape, accepts the sender's input to the OT functionality, and terminates.  

Next, we give the simulation of the receiver.  The simulator submits input to $\F_{\oslone}$, and receives output $u$.  Upon receiving OT
choice from the adversary, it feeds $u$ to the adversary as the OT
output. 
The simulation is perfect, since the view of the adversary contains nothing more than 
$u$. 
\end{proof}

